\documentclass{article}[12pt]
\usepackage{enumitem}
\usepackage{hyperref}
\usepackage{xcolor}
\usepackage[T1]{fontenc} % Should be Times New Roman adjacent, from https://tex.stackexchange.com/questions/33677/setting-font-family-for-the-whole-document
\usepackage{mathptmx}
\usepackage{moresize}
\usepackage{setspace}
\usepackage{float} 
\usepackage{placeins}
\usepackage{amsmath}
\usepackage{amssymb}
\usepackage{mathtools}
\usepackage[separate-uncertainty = true]{siunitx}
\usepackage{graphicx}
\usepackage{multicol}
\usepackage[margin=2.5cm]{geometry}
\usepackage{listings}

\graphicspath{{./images/}}
\pdfoutput=1
\onehalfspacing

\nonstopmode

\hypersetup{
    colorlinks,
    linkcolor={black!100!black},
    urlcolor={black!100!black}
}

\sisetup{
    number-unit-product=~
    } % Taken from M\X's answer on https://tex.stackexchange.com/questions/160105/unit-spacing-with-siunitx-in-an-equation

\newcommand{\uhref}[2]{\href{#1}{\underline{#2}}} % For directly linking to websites.
\newcommand{\reference}[3]{(#2, #3)\textsuperscript{\ref{#1}}} % For doing Harvard style referencing. Label should go at the bottom, reference at value being referenced.
\newcommand{\thlineA}{\hline ~ & ~ & ~ & ~ & ~ \\[-2ex] } % The values in the table get squished when using the regular hline, this adds a small bit of extra space at the top by adding a phantom row with negative bottom spacing.
\newcommand{\thline}{\hline ~ & ~ & ~ & ~ \\[-2ex] }


\begin{document}
    \section*{Abstract}
    \section{Introduction}
        \subsection{Background}


        \subsection{Theory}
            

    \section{Experiment}
        \subsection{Apparatus}
            \begin{itemize}

            \end{itemize}

        \subsection{Method}

            \subsubsection{Part 1}


            \subsubsection{Part 2}

    \section{Results}
        \begin{table}[!ht]
            \centering
            \begin{tabular}[center]{|c|c|c|c|}

            \end{tabular}
            \caption{Constants and accepted values}\label{constants}
        \end{table}
        \subsection{Part 1}

        \subsection{Part 2}
            
        \subsection{Error Analysis}
            \subsubsection{Quantitative}
                
            \subsubsection{Qualitative}
            
    \section{Discussion}
       
    \section{Conclusion}
        
    \addcontentsline{toc}{section}{Notes}
    \section*{Notes}

    \addcontentsline{toc}{section}{References}
    \section*{References}
    %https://tex.stackexchange.com/questions/42905/enumerated-list-with-square-brackets
        \begin{enumerate}[label={[\arabic*]}]
            \item
        \end{enumerate}
\end{document}